\documentclass[10pt,a4paper]{amsart}
\usepackage[margin=19mm]{geometry}
\usepackage{amsmath,amssymb,mathtools}
\usepackage[scr=rsfs,cal=euler]{mathalfa}
\usepackage{xfrac}
\usepackage[unicode,hidelinks,bookmarks=false]{hyperref}
\newcommand{\ie}{i.e.\ }
\newcommand{\cf}{cf.\ }
\newcommand{\resp}{respectively\ }
\newcommand{\Subsection}{\subsection}
\providecommand{\nobreakdash}{\nobreak\hbox{-}}
\setlength{\emergencystretch}{2em}
\multlinegap0pt
\usepackage{amssymb}
\usepackage{tikz}
\usepackage{float}
\newtheorem{proposition}{Proposition}
\newtheorem{theorem}{Theorem}
\newcommand{\BW}{\mathrm{BW}} 
\newcommand{\Fin}{\mathrm{Fin}}
\newcommand{\I}{\cI}
\newcommand{\J}{\cJ}
\newcommand{\cI}{\mathcal{I}}
\newcommand{\cJ}{\mathcal{J}}
\newcommand{\cc}{\mathfrak{c}}
\newcommand{\fin}{\mathrm{Fin}}
\title{A new order for ideal sequential compactness}
\author{Adam Kwela, Dorota Lesner}
\date{Source: arXiv:2603.01114v2 (CC BY 4.0)}
\begin{document}
\maketitle

\noindent\emph{Source and licence.} A new order for ideal sequential compactness. Version arXiv:2603.01114v2 (CC BY 4.0), distributed by arXiv under the Creative Commons Attribution 4.0 International licence, http://creativecommons.org/licenses/by/4.0/, as recorded in the arXiv licence metadata for this exact version and shown on the abstract page https://arxiv.org/abs/2603.01114v2. Full licence notice: \texttt{licenses/arxiv-2603.01114-CC-BY-4.0.txt}.\par\smallskip

\noindent\textit{Editorial scope.} Counted unit: the article's proposition prop:Fin\_vs\_tall (source0.tex lines 401--408). The article's complete original proof of it is delivered with it (source0.tex lines 410--416, one \texttt{proof} environment of 7 lines). No mathematics is written, repaired or re-derived: the statement and the proof are the article's own lines, in the article's own notation, and no retained line is substituted or re-flowed.  Retained uncounted as the source-local context the proof needs: lines 335--350; lines 390--393.  Nothing was omitted from any retained span, so the package carries no omission record.  No retained line contains a citation, so the wrapper carries no bibliography and no bibliography entry.  No retained line contains a figure environment, an image-inclusion command or drawing code, so no image is delivered and the package declares no asset dependency.  Editorial formatting only, in addition: an AMS \texttt{amsart} preamble that restates the article's own declarations and macros used by the retained lines, namely the environments \texttt{proposition}, \texttt{theorem} on separate counters, together with the standard packages those lines need.  The retained environments print in this document as Proposition 1, Theorem 1 in order of appearance, whereas the article prints the same items with the numbering of its own full document.  The complete native source of the article as distributed in the arXiv e-print for this exact version is bundled unchanged as \texttt{sources/195541/source0.tex}.
\begin{proposition}
\label{prop:basic-properties}
Let $\I$, $\J$ and $\mathcal{K}$ be ideals.
\begin{itemize}
    \item[(a)] The relation $\leq_{BW}$ is reflexive and transitive.
    \item[(b)] $\I\oplus\J\leq_{BW}\I$.
    \item[(c)] $\I\leq_{BW}\I \otimes \J$.
    \item[(d)] If $\I\leq_{BW}\J$, then $\I\leq_{K}\J$.
    \item[(e)] There is a $\leq_{BW}$-antichain of cardinality $\cc$ consisting of summable ideals.
    \item[(f)] The relation $\leq_{BW}$ is invariant over isomorphisms of ideals, i.e. if $\I'$ is isomorphic to $\I$ and $\J'$ is isomorphic to $\J$, then
    \[
    \I\leq_{BW}\J\ \iff\ \I'\leq_{BW}\J'.
    \]
    \item[(g)] If $\mathcal{K}\leq_{BW}\I$ and $\mathcal{K}\leq_{BW}\J$ then $\mathcal{K}\leq_{BW}\I \oplus \J$.
\end{itemize}
\end{proposition}

\begin{theorem}
\label{thm:zawieranieBW}
Let $\I$ and $\J$ be ideals. If $\I\leq_{BW}\J$, then $\BW(\J)\subseteq\BW(\I)$.
\end{theorem}

\begin{proposition}
\label{prop:Fin_vs_tall}
Let $\I$ be an ideal.
\begin{itemize}
    \item[(a)] $\I\leq_{BW}\Fin$ if and only if $\I$ is not tall. In particular, if $\I$ is not tall, then $\BW(\I)$ contains all sequentially compact spaces. 
    \item[(b)] If $\I$ is nowhere tall, then $\I$ and $\Fin$ are $\leq_{BW}$-equivalent. In particular, if $\I$ is nowhere tall, then $\BW(\I)$ is the class of all sequentially compact spaces.
\end{itemize}
\end{proposition}

\begin{proof}
(a): The implication $\impliedby$ follows from Proposition \ref{prop:basic-properties}(b), since any non-tall ideal is isomorphic to $\Fin\oplus\J$ for some ideal $\J$. On the other hand, the implication $\implies$ follows from Proposition \ref{prop:basic-properties}(d), since $\I\not\leq_K\Fin$ for tall ideals $\I$. The second part follows from Theorem \ref{thm:zawieranieBW} and the fact that $\BW(\Fin)$ is the class of all sequentially compact spaces.

(b): Let $\I$ be nowhere tall. Then $\I\leq_{BW}\Fin$ by item (a). We will show that $\Fin\leq_{BW}\I$. The second part will then follow from Theorem \ref{thm:zawieranieBW} and the fact that $\BW(\Fin)$ is the class of all sequentially compact spaces.

Without loss of generality we may assume that $\I$ is an ideal on $\omega$. Let $f:\omega \to \omega$ be the identity function. Fix $A \in \I^+$. Since $\I$ is nowhere tall, we can find an infinite $B\subseteq A$ such that $\I|B=\Fin|B$. Then $B = f[B]\subseteq f[A]$ and $B \in \fin^+$. Fix $C$ such that $C\cap B \in \fin^+$ and observe that then $A \cap f^{-1}[C]=A \cap C \supseteq B\cap C\in \I^+$.
\end{proof}

\end{document}
